\documentclass[10pt,a4paper]{article}

\usepackage[margin=1.15cm, top=1.1cm, bottom=1.2cm]{geometry}
\usepackage{fontspec}
\usepackage[russian,english]{babel}
\setmainfont{Liberation Serif}
\setsansfont{Liberation Sans}
\setmonofont{Liberation Mono}[Scale=0.83]

\usepackage{microtype}
\usepackage{booktabs}
\usepackage{tabularx}
\usepackage{array}
\usepackage{enumitem}
\usepackage[table]{xcolor}
\usepackage{titlesec}
\usepackage{hyperref}

\definecolor{HeaderBg}{HTML}{1E3A5F}
\definecolor{RowAlt}{HTML}{F4F7FA}
\definecolor{GreenTag}{HTML}{1B5E20}
\definecolor{AmberTag}{HTML}{B26A00}
\definecolor{RedTag}{HTML}{B71C1C}

\hypersetup{
    colorlinks=true,
    linkcolor=HeaderBg,
    urlcolor=HeaderBg,
    pdftitle={Анализ стратегий чанкинга и Hybrid Search — Навигатор по регламентам МФТИ}
}

\titleformat{\section}{\normalsize\bfseries\sffamily\color{HeaderBg}}{\thesection.}{0.4em}{}[\titlerule]
\titleformat{\subsection}{\small\bfseries\sffamily\color{HeaderBg}}{\thesubsection.}{0.35em}{}
\titlespacing*{\section}{0pt}{5pt}{2pt}
\titlespacing*{\subsection}{0pt}{3.5pt}{1.5pt}

\setlist[itemize]{noitemsep, topsep=1pt, parsep=0.5pt, leftmargin=1.1em}
\setlist[enumerate]{noitemsep, topsep=1pt, parsep=0.5pt, leftmargin=1.3em}
\setlength{\parindent}{0pt}
\setlength{\parskip}{1.5pt}
\setlength{\emergencystretch}{2em}
\renewcommand{\arraystretch}{1.07}

\newcommand{\keep}{\textbf{\textcolor{GreenTag}{[В ЯДРО]}}}
\newcommand{\modify}{\textbf{\textcolor{AmberTag}{[ДОРАБОТАТЬ]}}}
\newcommand{\reject}{\textbf{\textcolor{RedTag}{[ОТКЛОНИТЬ]}}}

\begin{document}
\small
\begin{center}
    {\Large\bfseries\sffamily\color{HeaderBg} Аудит идей по чанкингу и Hybrid Search: «Навигатор по регламентам МФТИ»}\\[2pt]
    {\small Сравнительная оценка материалов \texttt{ответ\_соннет.txt} (Д1) и \texttt{Чанкинг и Hybrid Search - анализ (2).docx} (Д2) на основе верификации корпуса из 20 файлов и критериев оценки ТЗ}
\end{center}
\vspace{-6pt}

\section{Привязка к ТЗ и верифицированной структуре датасета}
Архитектура парсинга, чанкинга и поиска напрямую определяет \textbf{52 из 100 технических баллов} жюри:
\textbf{К1 (20 б.)} — извлечение структуры вплоть до пунктов/подпунктов (8 б.), привязка фрагментов к документу/разделу/пункту (6 б.), масштабируемость без ручной правки кода (5 б.);
\textbf{К2 (10 б.)} — учёт дат/версий (4 б.) и точное указание редакции и пункта (6 б.);
\textbf{К3 (22 б.)} — релевантность поиска (6 б.), полнота извлечения шагов процедуры (8 б.), учёт ролей и условий (5 б.);
дополнительно в \textbf{К4 (20 б.)} — явный показ противоречий между документами (3 б.) и отказ при отсутствии ответа в базе (5 б.).

\textbf{Фактическая ревизия папки \texttt{данные/} (20 файлов: 9 PDF на 79 стр., 8 DOCX, 3 XLSX):}
\begin{itemize}
    \item \textbf{Сканы без текстового слоя — 5 из 9 PDF (44 из 79 стр.):} \texttt{460-1 Открытие ФЛС.pdf} (19 стр., рукописный № от 28.04.2017), \texttt{Регламент бронирования конц зала.pdf} (14 стр., штамп № 3897-1 от 01.09.2025), \texttt{Приказ и регламент по ЭБ.pdf} (6 стр., № 1027-1 от 25.06.2019, \emph{таблицы на стр. 5--6 повернуты на 90°}), \texttt{Регламент эксплуатации ЦОД.pdf} (5 стр., № 1292-1 от 28.03.2025), \texttt{Регламент поселения сотрудников.pdf} (3 стр., 2024 г.; \emph{дефект исходного скана: стр.~1 обрывается на заголовке Раздела~2, а стр.~2 начинается сразу с п.~2.6 — пункты 2.1--2.5 физически отсутствуют}).
    \item \textbf{Версионные цепочки и коллизии:} (1) \textbf{ФЛС:} приказ \texttt{460-1} (2017, бумажный обход через ПФО) $\to$ приказ \texttt{3065-1} (25.11.2021, изменяет \texttt{1385-1} от 2020, переводит открытие ФЛС в \texttt{1С:ДО} и вводит учёт РИД) + распоряжение \texttt{700} (28.12.2024, 15-значная структура кода ФЛС). (2) \textbf{Сметы:} приказы \texttt{794-1} (28.03.2023) и \texttt{889-1} (04.04.2023) + бланки 2023 г. (\texttt{Приложение 1а/1б}) против старой формы 2019 г. (\texttt{Приложение 1\_новая смета 2019.xlsx}). (3) \textbf{Суточные:} приказ \texttt{3699-1} (18.08.2025, п.~17: 100 руб. работникам, 50 руб. обучающимся) против проекта \texttt{Регламент о командировании.pdf} (2025, без номера: «не более 100 руб.»). (4) \textbf{Заимствования:} приказ \texttt{5990-1} (03.11.2026) отменяет приказ \texttt{1465-1} (2020).
    \item \textbf{Аномалии вёрстки:} в 8 DOCX нет автонумерации \texttt{w:numPr} (все номера — обычный текст) и сломаны стили \texttt{Heading}; внутри одного файла до 6 независимых единиц (\texttt{Приказ} + \texttt{Приложения 1..N}) со сбросом нумерации с \texttt{1.}; в \texttt{Приказ N 3065-1} дублируются номера пунктов \texttt{2.6} и \texttt{2.7}; в DOCX лежат огромные таблицы \texttt{ЛИСТ СОГЛАСОВАНИЯ}.
\end{itemize}

\section{Сводная матрица оценки предложенных идей}

{\footnotesize
\rowcolors{2}{RowAlt}{white}
\begin{tabularx}{\textwidth}{@{} >{\raggedright\arraybackslash}p{0.27\textwidth} c >{\raggedright\arraybackslash}p{0.135\textwidth} >{\raggedright\arraybackslash}X @{}}
\toprule
\rowcolor{HeaderBg}
\textcolor{white}{\textbf{Полезные и дорабатываемые идеи}} & \textcolor{white}{\textbf{Ист.}} & \textcolor{white}{\textbf{Вердикт}} & \textcolor{white}{\textbf{Краткое обоснование относительно проекта}} \\
\midrule
Дерево \texttt{file $\to$ unit $\to$ section $\to$ clause $\to$ item} + \textbf{V3 (Parent-Child)} & Д1, Д2 & \keep & Закрывает К1 (14 б.) и К3 (8 б.): индекс по \texttt{clause} ($\le 1200$ симв.) даёт точный пункт, подтягивание \texttt{section} даёт полный алгоритм шагов. \\
Разделение синтетического \texttt{chunk\_id} и отображаемого \texttt{clause\_no} & Д1, Д2 & \keep & Критично из-за сброса нумерации в приложениях и дублей номеров пунктов (\texttt{2.6}, \texttt{2.7} в \texttt{3065-1}) внутри одного файла. \\
Детерминированный префикс (\textit{хлебные крошки}) в \texttt{embed\_text} & Д1, Д2 & \keep & Без префикса \texttt{[Приказ · Приложение · Раздел]} короткие пункты («2.3. Срок — 3 р.д.») теряют смысл при векторизации. \\
Сегментация зон и исключение \texttt{ЛИСТ СОГЛАСОВАНИЯ} из поиска & Д1, Д2 & \keep & Таблицы согласования 1С:ДО содержат десятки ФИО и должностей, замусоривая выдачу по ролевым запросам (К3). \\
Лемматизированный BM25 (\texttt{pymorphy3}) + глоссарий + Dense (\texttt{BGE-M3}) + \texttt{exact=True} & Д2 (Д1) & \keep & На $\sim$800 чанках полный перебор \texttt{exact=True} в Qdrant занимает $<5$ мс; лемматизация и расшифровка (\texttt{ФЛС}, \texttt{КФО}, \texttt{ЦФО}) обязательны для русского языка. \\
Парсинг DOCX через \texttt{iterchildren()} + дедупликация слитых ячеек & Д2 & \keep & Стандартный обход \texttt{doc.paragraphs} теряет позицию таблиц внутри пунктов, а объединённые ячейки дублируют текст в \texttt{python-docx}. \\
Forward-fill иерархии в XLSX + чанки по 10--15 строк с шапкой + формулы смет & Д1, Д2 & \keep & В \texttt{Фин. структуре 2026} (525 строк) родительские ЦФО указаны только в верхней ячейке группы; в сметах 2019/2023 правила зашиты в формулах (\texttt{0.082}, \texttt{0.2}). \\
Калиброванный порог реранкера для детекции «ответа нет в базе» & Д1, Д2 & \keep & Прямое требование К4 (5 б.): отказ должен срабатывать до генерации или через строгий флаг при \texttt{max\_rerank\_score < $\tau$}. \\
Чек-лист приёмки парсера + золотой датасет на 50 вопросов & Д2 & \keep & Единственный способ гарантировать отсутствие тихих потерь пунктов и регрессий перед защитой. \\
\midrule
Реестр версий (\texttt{supersedes}, \texttt{status}) + YAML-оверрайд метаданных & Д1, Д2 & \modify & Вести граф версий обязательно, но YAML должен быть \emph{кэшем/оверрайдом поверх автоэкстрактора}, чтобы не нарушить К1 (5 б. за новые файлы без правки кода). \\
Фильтрация \texttt{status != superseded} в Qdrant по умолчанию & Д1, Д2 & \modify & \textbf{Нельзя резать жёстким фильтром!} Иначе система потеряет 5 приложений базового приказа \texttt{460-1} и не сможет показать противоречия редакций (К2, К4). Нужен мягкий бустинг + парная подгрузка версий. \\
Дедупликация парных файлов (\texttt{3614-1} и \texttt{3699-1}) & Д1, Д2 & \modify & \texttt{sha256(file)} из Д1 не сработает (разные форматы DOCX и XLSX, разный состав). Дедуплицировать нужно на уровне \texttt{unit}/таблицы с сохранением обоих \texttt{doc\_id}. \\
\bottomrule
\end{tabularx}

\vspace{4pt}
\rowcolors{2}{RowAlt}{white}
\begin{tabularx}{\textwidth}{@{} >{\raggedright\arraybackslash}p{0.27\textwidth} c >{\raggedright\arraybackslash}p{0.135\textwidth} >{\raggedright\arraybackslash}X @{}}
\toprule
\rowcolor{HeaderBg}
\textcolor{white}{\textbf{Отклоняемые идеи и ошибки}} & \textcolor{white}{\textbf{Ист.}} & \textcolor{white}{\textbf{Вердикт}} & \textcolor{white}{\textbf{Причина отказа / почему вредит проекту}} \\
\midrule
Тезис об отсутствии 3 регламентов (поселение, подбор, ЦОД) и всего 3 сканах & Д1 & \reject & \textbf{Фактическая ошибка Д1:} все 20 файлов есть в папке \texttt{данные/}, а сканов 5, а не 3. Отказ отвечать по этим темам обнулит баллы на базовых вопросах. \\
Ставка на голый \texttt{Tesseract OCR} как основной движок для сканов & Д1, Д2 & \reject & В системе нет \texttt{rus.traineddata}; Tesseract даёт битые слова, не читает рукописные номера приказов (\texttt{460-1}, \texttt{1027-1}) и ломается на повёрнутых на 90° таблицах ЭБ. Нужен VLM-OCR с кэшированием по хэшу. \\
Двухэтапный поиск: сначала Top-K по «карточке документа», затем поиск пунктов & Д1 & \reject & На корпусе в 800 чанков это искусственное бутылочное горлышко: ошибка на 1-м шаге убивает полноту в мультидокументных вопросах (\texttt{3699-1} + \texttt{Командировки}). \\
LLM Contextual Retrieval (\textbf{V4}: 800 вызовов LLM при индексации) и \textbf{HyDE} & Д1, Д2 & \reject & Детерминированный шаблон хлебных крошек даёт 95\% эффекта бесплатно; 800 вызовов LLM при старте Docker-контейнера замедлят развёртывание (риск по К6 — 20 б.). \\
Fixed-window (\textbf{V0}), разрез только по разделам (\textbf{V1}), Semantic (\textbf{V5}), GraphRAG & Д1, Д2 & \reject & \textbf{V0/V5} разрушают границы пунктов (\texttt{п. 2.3.1}), лишая до 22 баллов за точное цитирование; \textbf{V1} даёт чанки по 18--26 тыс. символов (переполнение и шум). \\
\bottomrule
\end{tabularx}
}

\section{Детальный разбор полезных идей (что берём в реализацию и почему)}

\subsection{Пятиуровневое дерево документа и стратегия чанкинга V3 (Parent-Child)}
Оба документа (особенно Д2 в §3 и §5) справедливо отвергают плоскую нарезку текста и предлагают модель \texttt{file $\to$ unit $\to$ section $\to$ clause $\to$ item}.
\begin{itemize}
    \item \textbf{Почему уровень \texttt{unit} критичен:} в файлах МФТИ один физический документ содержит несколько юридически разных сущностей. Например, \texttt{Приказ N 3699-1.docx} включает сам Приказ (24 пункта), Лист согласования и 4 Приложения (Регламент формирования бюджета, План-график, Таблицу ЦФО). В каждом приложении нумерация разделов и пунктов начинается заново с \texttt{1.}. Без выделения \texttt{unit\_id} ссылка «п. 2.1» неоднозначна.
    \item \textbf{Почему именно V3 (Parent-Child):} индексация на уровне \texttt{clause} (пункт \texttt{2.3} вместе с его подпунктами \texttt{а, б, в}, если суммарная длина $\le 1200$ символов) обеспечивает снайперское попадание в конкретную норму (К1: 8 б., К2: 6 б.). При этом хранение ссылки \texttt{parent\_id} на весь родительский \texttt{section} (и ссылок \texttt{prev\_id}/\texttt{next\_id}) позволяет на этапе сборки контекста перед LLM подтянуть соседние пункты раздела и выдать полный пошаговый план действий без пропусков (К3: 8 б.).
    \item \textbf{Синтетический \texttt{chunk\_id} $\neq$ \texttt{clause\_no}:} в \texttt{Приказ N 3065-1.docx} из-за опечатки авторов регламента дважды встречаются пункты \texttt{2.6} и \texttt{2.7} с разным текстом. Использование порядкового синтетического ключа (\texttt{3065-1::reg::s2::c08}) при сохранении оригинального \texttt{clause\_no="2.6"} для цитирования предотвращает перезапись чанков в БД.
\end{itemize}

\subsection{Детерминированное контекстное обогащение (\texttt{embed\_text} $\neq$ \texttt{raw\_text})}
В Qdrant сохраняется чистый \texttt{raw\_text} (для цитирования и показа пользователю), а эмбеддинг и лемматизированный BM25 строятся по полю \texttt{embed\_text}, снабжённому префиксом:
\texttt{[Документ: Приказ № 3065-1 от 25.11.2021 (ред. приказа 1385-1) · Статус: действует · Единица: Регламент открытия ФЛС · Раздел 3: Порядок открытия ФЛС · Пункт 3.4 · Роли: Руководитель проекта, ПФО, Бухгалтерия]} + текст пункта.
\textbf{Аргумент:} свыше 40\% пунктов в регламентах МФТИ не содержат названия процедуры внутри самого абзаца (например: \emph{«В течение 3 рабочих дней направляет карточку на согласование»}). Префикс связывает такие пункты с темой документа без платных вызовов LLM.

\subsection{Специализированная обработка таблиц, Excel-форм и листов согласования}
\begin{itemize}
    \item \textbf{Отсечение зоны \texttt{approval}:} в 8 DOCX-приказах в конце документа расположены таблицы \texttt{ЛИСТ СОГЛАСОВАНИЯ} из \texttt{1С:ДО} с фамилиями проректоров и начальников управлений. Если их индексировать наравне с текстом, запрос \emph{«кто согласует смету проекта?»} вернёт список подписантов самого приказа вместо ролей по процедуре. Разметка \texttt{zone="approval"} и исключение её из основного поиска решают проблему.
    \item \textbf{Forward-fill и нарезка XLSX/DOCX-таблиц:} в \texttt{Финансовая структура 2026} (525 строк) коды ЦФО 1-го и 2-го уровня объединены вертикально (\texttt{merged\_cells}). При построчном чтении дочерние лаборатории теряют привязку к физтех-школе/институту. Протяжка значений сверху вниз (\texttt{ffill}) и упаковка по 10--15 строк одного подразделения с повтором шапки таблицы сохраняют табличный контекст.
    \item \textbf{Извлечение формул из бланков смет:} файлы \texttt{Приложение 1а/1б} и \texttt{новая смета 2019.xlsx} — это пустые шаблоны с 848 формулами. Нормативы накладных расходов (\texttt{=SUM(...)*0.2} для РЦНИ или \texttt{0.082}) и лимиты ФОТ зашиты прямо в формулы ячеек. Парсер должен генерировать текстовую карточку структуры бланка и коэффициентов из формул.
\end{itemize}

\subsection{Гибридный поиск в Qdrant: \texttt{exact=True}, лемматизированный BM25 и глоссарий}
\begin{itemize}
    \item \textbf{Точный поиск (\texttt{exact=True}):} объём всего корпуса после чанкинга составляет $\sim$750--900 чанков. На таком масштабе приближённый граф HNSW не даёт выигрыша в скорости, но может терять редкие точные совпадения при фильтрации. Полный перебор (\texttt{exact=True}) работает за единицы миллисекунд со 100\% полнотой.
    \item \textbf{Клиентский Sparse BM25 с \texttt{pymorphy3} + словарь аббревиатур:} стандартный токенизатор без русской морфологии не связывает формы \emph{«командирование» / «командировке» / «командированный»}. Кроме того, в регламентах есть разделы «Термины и сокращения» (\texttt{ФЛС}, \texttt{КФО}, \texttt{ЦФО}, \texttt{РИД}, \texttt{УИТ}, \texttt{ФЭУ}, \texttt{ОКР}). Автоматический сбор пар \texttt{аббревиатура $\leftrightarrow$ расшифровка} при парсинге и двустороннее расширение текста перед BM25 позволяют находить документ как по аббревиатуре, так и по полному названию.
\end{itemize}

\section{Аргументированный разбор ошибочных и вредных идей (что отклоняем)}

\subsection{Фактическая ошибка в \texttt{ответ\_соннет.txt} о «недостающих регламентах» и числе сканов}
В \texttt{ответ\_соннет.txt} (раздел 1 и п. 2.2) утверждается, что в датасете якобы всего 3 скана, а \emph{«Регламент поселения в общежитие, Регламент поиска и подбора персонала, Регламент эксплуатации ЦОД отсутствуют в датасете»}, и предлагается настроить систему на отказ по этим вопросам.
\begin{itemize}
    \item \textbf{Почему это грубая ошибка:} в папке \texttt{данные/} лежат ровно 20 файлов, включая \texttt{Регламент поселения сотрудников.pdf} (скан, 3 стр.), \texttt{Регламент эксплуатации центра обработки данных МФТИ, Физтех.pdf} (скан, 5 стр.) и \texttt{Регламент поиска и подбора персонала.pdf} (текстовый PDF, 10 стр.). Автор Д1 пропустил 2 файла при обходе директории и не распознал сканы. Настройка отказа по этим темам приведёт к потере баллов на прямых вопросах жюри.
    \item \textbf{Реальный нюанс с поселением, который нужно обрабатывать:} в файле \texttt{Регламент поселения сотрудников.pdf} действительно есть дефект сканирования — между стр.~1 и стр.~2 пропущены пункты \texttt{2.1--2.5} (текст перескакивает с заголовка раздела 2 сразу на \texttt{п. 2.6}). Вот об этой неполноте исходного скана система должна честно предупреждать пользователя!
\end{itemize}

\subsection{Опасность жёсткого фильтра \texttt{status == active} (\texttt{status != утратил}) по умолчанию}
Оба документа предлагают по умолчанию фильтровать неактуальные документы на этапе векторного поиска в Qdrant.
\begin{itemize}
    \item \textbf{Почему это навредит на хакатоне:}
    Во-первых, критерии \textbf{К2 (10 б.)} и \textbf{К4 (3 б.)} требуют, чтобы система распознавала разные редакции и \textbf{явно показывала пользователю расхождения и изменения между версиями}. Если чанки приказа \texttt{460-1} (2017) отсечены жёстким фильтром до поиска, LLM даже не узнает о существовании старой редакции и не сможет объяснить, что бумажный порядок через ПФО заменён на электронный в \texttt{1С:ДО} приказом \texttt{3065-1} (2021).
    Во-вторых, изменяющие приказы часто не дублируют все приложения базового приказа: например, в \texttt{460-1} содержится 5 приложений (включая Перечень критических технологий и шаблоны распоряжений), а в \texttt{3065-1} — только обновлённый текст Регламента и Карточка проекта (Приложения 1 и 2). Жёсткий фильтр \texttt{status != superseded} удалит из базы действующие приложения базового приказа.
    \item \textbf{Правильное решение (Version-Aware Expansion):} искать по всем версиям (кроме зоны \texttt{approval}), применять мягкий бустинг к действующим редакциям на этапе реранкинга, а при попадании в выдачу чанка со связью \texttt{supersedes} / \texttt{superseded\_by} автоматически подгружать парный пункт другой редакции и передавать оба в LLM с явной меткой версий.
\end{itemize}

\subsection{Неприменимость голого \texttt{Tesseract OCR} как основного инструмента для 5 сканов}
Оба автора предлагают использовать \texttt{pdf2image + Tesseract (rus+eng)}.
\begin{itemize}
    \item \textbf{Почему это плохое решение:}
    (1) На целевой машине в \texttt{/usr/share/tessdata} вообще не установлен пакет \texttt{rus.traineddata}.
    (2) Как видно из самого файла \texttt{ответ\_соннет.txt}, Tesseract даёт критический уровень шума на наших сканах (\emph{«оборулования»}, \emph{«ст 'уденческие мероприя Гия»}), из-за чего не срабатывают регулярные выражения для номеров пунктов и таблиц.
    (3) В \texttt{460-1} и \texttt{Приказ и регламент по ЭБ.pdf} номера и даты приказов вписаны \textbf{от руки} в штамп регистрации — Tesseract их не берёт.
    (4) В \texttt{Приказ и регламент по ЭБ.pdf} на стр. 5--6 таблицы классификации помещений \textbf{повёрнуты на 90° (альбомная ориентация внутри книжного скана)} — без детекции ориентации или мультимодального чтения таблица превращается в кашу символов.
    \item \textbf{Правильное решение:} использовать Vision-LLM (VLM OCR в Markdown с сохранением таблиц и чтением рукописных штампов) для 44 страниц сканов с \textbf{кэшированием результата на диск по \texttt{sha256(file)}} (чтобы при перезапуске контейнера индексация занимала $<2$ сек), оставив связку \texttt{Docling / RapidOCR / Tesseract} только как офлайн-фоллбэк для новых файлов.
\end{itemize}

\subsection{Отказ от двухэтапного поиска по «карточке документа» (\texttt{ответ\_соннет.txt}, п. 3.2)}
Автор Д1 предлагает сначала искать Top-2 документа по сгенерированным саммари («карточкам документов»), и только внутри них искать пункты.
\begin{itemize}
    \item \textbf{Почему это антипаттерн для нашего корпуса:} во всём датасете всего 20 документов и $\sim$800 чанков. Двухэтапная фильтрация создаёт узкое горлышко: если вопрос требует синтеза нескольких документов (например, \emph{«какие нормы суточных и как их заложить в смету на 2026 год?»} $\to$ нужны \texttt{Приказ 3699-1}, \texttt{Регламент о командировании}, \texttt{Приказ 794-1} и \texttt{Распоряжение 700}), отсечение на первом шаге потеряет половину источников и обрушит баллы по К3 и К4. Глобальный поиск по всем 800 чанкам с контекстным префиксом одновременно учитывает и тему документа, и текст конкретного пункта без потери полноты.
\end{itemize}

\subsection{Отказ от дедупликации по хэшу файла (\texttt{sha256(file)})}
В \texttt{ответ\_соннет.txt} (п. 1.2) предлагается удалять дубликаты файлов в \texttt{ingest.py} по хэшу файла.
\begin{itemize}
    \item \textbf{Почему не сработает:} дублирующиеся приложения в датасете имеют \textbf{разный формат и бинарный состав}: \texttt{Приказ N 3614-1.docx} включает текст приказа + таблицу финструктуры в Word, а \texttt{Приложение к приказу...xlsx} — ту же таблицу в Excel. Аналогично, \texttt{Приказ N 3699-1.docx} содержит приказ + регламент + приложения, а файл \texttt{Приложения к приказу\_ о формировании бюджетирования 2026г..docx} — только приложения. Их \texttt{sha256} различны. Вместо удаления файла нужно связывать их в реестре как \texttt{master\_doc} $\leftrightarrow$ \texttt{standalone\_appendix}, индексируя табличную структуру из XLSX (где сохранена иерархия ячеек), но указывая в метаданных оба источника для корректного цитирования.
\end{itemize}

\section{Итоговая целевая архитектура пайплайна (Synthesis)}
\begin{enumerate}
    \item \textbf{Ingestion \& Normalization:} Маршрутизатор по типу файла (\texttt{DOCX} через \texttt{body.iterchildren()} с очисткой слитых ячеек и отсечением \texttt{ЛИСТ СОГЛАСОВАНИЯ}; \texttt{PDF-text} через \texttt{PyMuPDF/pdfplumber} со склейкой переносов и очисткой сносок вида \texttt{8123 $\to$ 812}; \texttt{PDF-scan} через кэшированный \texttt{VLM-OCR} с поддержкой повёрнутых таблиц и рукописных штампов; \texttt{XLSX} через \texttt{openpyxl} с \texttt{forward-fill} объединённых ячеек и экстракцией формул смет).
    \item \textbf{Document Registry \& Version Graph:} Автоматический парсер шапки приказа (\texttt{doc\_number}, \texttt{doc\_date}, триггеры \emph{«Во изменение приказа...»}, \emph{«признать утратившим силу...»}) + верифицированный \texttt{registry\_overrides.yaml} для 20 эталонных файлов.
    \item \textbf{Hierarchical Chunker (V3 Parent-Child):} Разбиение на \texttt{unit $\to$ section $\to$ clause $\to$ item}. Векторизация на уровне \texttt{clause} ($\le 1200$ симв.) и табличных групп (10--15 строк с шапкой) с детерминированным префиксом \texttt{[Документ · Статус · Раздел · Пункт · Роли]}.
    \item \textbf{Hybrid Search + Reranking + Context Assembly:} Qdrant (\texttt{exact=True}): \texttt{Dense (BGE-M3)} + \texttt{Sparse BM25 (pymorphy3 + глоссарий аббревиатур)} $\to$ \texttt{RRF (top 30)} $\to$ \texttt{Cross-Encoder Reranker (BGE-reranker-v2-m3, top 8)}. Проверка порога отказа (\texttt{max\_score < $\tau$}). Расширение контекста: подгрузка родительского \texttt{section}, соседних пунктов и связанных версий (\texttt{supersedes}/\texttt{superseded\_by}) для явного отображения коллизий в ответе.
\end{enumerate}

\end{document}
